\section{Algorithms}\label{app:algorithms}

This appendix specifies every computation used in the paper, in enough detail to be
reimplemented independently. Two kinds of computation occur.
\begin{itemize}
\item \emph{Exact} computations (\cref{alg:frontier,alg:lift,alg:glue,alg:moves,alg:certificate,alg:seeds,alg:bregman,alg:edges})
use integer arithmetic only. Their outputs either are certificates that can be checked
independently (\cref{app:data-cert,app:data-glue}) or support statements labelled
``computer-verified''.
\item \emph{Numerical} computations (\cref{alg:sinkhorn-cont,alg:bethe,alg:predictor}) use
double-precision floating point. They support only the numerical evidence for the
conjectures, never a theorem. We report convergence residuals with each result.
\end{itemize}
Throughout, $\mathrm{Adj}_n(v)=\{k^2-v:\ \lfloor\sqrt v\rfloor<k\le\lfloor\sqrt{v+n}\rfloor\}\setminus\{v\}$ is the
neighbourhood of $v$ in $G(n)$, and $\operatorname{issq}(m)$ tests whether $m$ is a perfect square
using an integer square root. In particular,
\begin{equation}\label{eq:deg}
\deg_{G(n)}(v)=\lfloor\sqrt{v+n}\rfloor-\lfloor\sqrt v\rfloor-[\,2v\text{ is a square}\,].
\end{equation}

\subsection{Exact counting of Hamiltonian paths}
The values of $\Ham(n)$ in \cref{tab:counts} were computed in two stages.

\emph{Graphillion stage.} For $n\le56$ we first used the library Graphillion \cite{graphillion},
which represents sets of subgraphs as zero-suppressed decision diagrams \cite{Minato1993} and
builds them by frontier-based search in the spirit of Knuth's \textsc{Simpath}
\cite[\S7.1.4]{KnuthTAOCP4A}. On the advice of S.~Minato, the vertex order was an optimal
\emph{vertex-separation} ordering, whose width equals the pathwidth \cite{Kinnersley}. It was
computed exactly by encoding Coudert's integer-programming formulation \cite{Coudert} as a weighted
MaxSAT instance. This approach exhausted the available storage at about $n=56$.

\emph{Dynamic programming stage.} We then switched to a dynamic program over the same ordering.
It is a Held--Karp-type recursion \cite{Bellman,HeldKarp} restricted to a path decomposition, so
that a state describes only the frontier rather than an arbitrary subset of vertices. This is the
approach of Pettersson \cite{Pettersson}, which we found only afterwards. The frontier after
step~$t$ is the set of processed vertices that still have an unprocessed neighbour. Its size is at
most the width $w$ of the ordering; for example $w=11$ for $n=57$ and $w=13$ for $n=59$.

\emph{States.} A state is a tuple $(\sigma_f)_{f\in F}$ together with a counter $c\in\{0,1\}$.
\begin{itemize}
\item $\sigma_f=0$ if $f$ has no chosen edges yet.
\item $\sigma_f=-1$ if $f$ has degree $2$, i.e.\ is interior to a chain.
\item $\sigma_f=L>0$ if $f$ is an endpoint of the partial path segment (``chain'') labelled $L$.
\item $c$ records whether a chain has been \emph{closed}, meaning both of its endpoints have
left the frontier.
\end{itemize}
Labels are canonicalised by renumbering them in order of first appearance, so there are at most
$\mathrm{Bell}(w)$ connectivity patterns per frontier. The counter $c$ is essential for
correctness. Without it, two half-external chains (each with one endpoint already gone) that are
joined by an edge leave no trace in $\sigma$, and the resulting phantom states are counted as
Hamiltonian paths. For $n=15$ the naive version returns $11$ instead of $1$.

\begin{algorithm}[h]
\caption{Frontier dynamic program for $\Ham(n)$ along a pathwidth ordering}\label{alg:frontier}
\begin{algorithmic}[1]
\Require $n$; an ordering $v_1,\dots,v_n$ of $[n]$ of small width (exact vertex separation, then refined)
\State $\mathcal T\gets\{(\emptyset,c{=}0)\mapsto1\}$ \Comment{state $\to$ number of partial edge sets}
\For{$t=1,\dots,n$}
  \State $N_t\gets$ neighbours of $v_t$ among the processed vertices (all of them are in the frontier)
  \State $\mathcal T'\gets\emptyset$
  \For{each $(\sigma,c)\mapsto\#$ in $\mathcal T$}
    \For{each $X\subseteq\{u\in N_t:\sigma_u\ne-1\}$ with $|X|\le2$} \Comment{edges from $v_t$ to the past}
      \If{$X=\{u,u'\}$ with $u,u'$ the two endpoints of one chain} \textbf{continue} \Comment{cycle}\EndIf
      \State $\sigma'\gets\sigma$ updated by adding the edges $\{v_t,u\}$, $u\in X$: a vertex with
      \State \hspace{1.2em} status $0$ becomes an endpoint of a new or extended chain, an endpoint becomes
      \State \hspace{1.2em} interior ($-1$), chains joined through $v_t$ are merged, and $v_t$ gets status $0$, $L$ or $-1$
      \State $\mathcal T'[(\sigma',c)]\mathrel{+}=\#$
    \EndFor
  \EndFor
  \For{each $f$ with no unprocessed neighbour left} \Comment{fused elimination}
    \State in every state: discard if $\sigma_f=0$; delete $f$ if $\sigma_f=-1$;
    \State \hspace{1.2em} if $\sigma_f=L$, then $f$ becomes a path end: delete $f$, and if the other end of chain $L$
    \State \hspace{1.2em} has already left, set $c\gets c+1$; discard if $c>1$, or if $c=1$ while
    \State \hspace{1.2em} $t<n$ or other chains remain
  \EndFor
  \State canonicalise labels, merge equal states; $\mathcal T\gets\mathcal T'$
\EndFor
\State \Return $\mathcal T[(\emptyset,c{=}1)]$ \Comment{each undirected Hamiltonian path counted once}
\end{algorithmic}
\end{algorithm}

\emph{Ordering and engineering.} Among orderings of optimal width, the running time still
varies considerably. We therefore refined the exact vertex-separation ordering by simulated
annealing on a cost function that predicts the size of the state tables from the frontier
profile. The implementation is in C. States are packed into $64$-bit keys with $128$-bit counts,
and the edge sweep for $v_t$ (all subsets $X$ at once) is fused with the eliminations that
follow it. At each step the main cost is merging a stream of $(\text{key},\text{count})$ pairs,
in which many source states map to the same key, into the next table. This is done as follows.
\begin{itemize}
\item The next table is an extendible hash table \cite{Fagin}, so it grows in proportion to the
actual output.
\item New pairs are first collected in a radix-sort buffer indexed by a hash prefix, and are
inserted bucket by bucket, so that each bucket is brought into cache once per flush rather than
once per insertion.
\item Buffer and bucket sizes are chosen from the cache sizes of the machine.
\item Freed slabs are returned to the operating system explicitly.
\end{itemize}
With these measures, $n=61$ (width $14$; peak $6.2\cdot10^8$ states in one step) runs in
$97$ minutes in under $2$\,GB of memory on a laptop (Apple M3 Pro). The largest cases, up to $n=75$, were run on a machine with $768$\,GB of memory and $64$ cores,
with a parallel sort--merge version of the same step. Each source state and each subset of
back-edges contributes independently, and only the summation of counts with equal keys couples
them. So $P$ workers process disjoint contiguous slices of the current table. Each worker writes
its contributions into a private bounded buffer, which it sorts by a hash of the key and
deduplicates into sorted runs, kept in memory or spilled to SSD. The runs are then merged with
deduplication by $P$ threads, each owning a disjoint interval of hash space chosen from sampled
splitters. Because each table is stored in hash order, a contiguous slice is in effect a random
sample of states, so static slices balance the very uneven per-state work: the slowest worker
was within $20\%$ of the mean on the dominant steps. The run for $n=75$ is the only one that
reaches frontier width $15$; its value was confirmed by two runs with different optimal
orderings, after an earlier value computed with a $15$-slot state encoding was withdrawn. Details of the implementation will appear separately \cite{MillerEngineering}.

\subsection{The lifting step and the glue}

\begin{algorithm}[h]
\caption{Lift (\cref{lem:segments,prop:lift})}\label{alg:lift}
\begin{algorithmic}[1]
\Require odd $t$, $T=t^2$, $h=(T-1)/2$; a Hamiltonian path $b=(b_1,\dots,b_n)$ of $G(n)$, $n$ odd, from $1$ to $3$;
a glue arrangement $\Gamma$ (a list of glue vertices and segment pairs $[x\ y]$)
\For{$c=-h,\dots,h$}
  \State $T_c\gets(Tb_p+(-1)^{p+1}c)_{p=1}^n$ \Comment{from $T+c$ to $3T+c$}
\EndFor
\State $P'\gets(\,)$; $\mathrm{offset}\gets\{\}$
\For{each item of $\Gamma$ in order}
  \If{the item is a glue vertex $g$} append $g$ to $P'$
  \Else{ (a pair $[x\ y]$)} let $c$ be the index with $\{x,y\}=\{T+c,3T+c\}$;
    \State \hspace{1.2em} record $\mathrm{offset}[c]\gets(|P'|,\ \text{reversed}:=[x=3T+c])$;
    \State \hspace{1.2em} append $T_c$, reversed if $x=3T+c$, to $P'$
  \EndIf
\EndFor
\State \Return $P'$ (a Hamiltonian path of $G(Tn+h)$ from $1$ to $3$) and $\mathrm{offset}$
\end{algorithmic}
\end{algorithm}

\begin{algorithm}[h]
\caption{Glue search for $(t,u,w,\text{parity of }n)$}\label{alg:glue}
\begin{algorithmic}[1]
\State $h\gets(t^2-1)/2$; for $c\in[-h,h]$ let the forced pair be $\{t^2u+c,\ t^2w+\varepsilon c\}$, with $\varepsilon=+1$ for $n$ odd and $-1$ for $n$ even
\State $V\gets\{1,\dots,h\}\cup\{\text{all endpoints}\}$; join $x\ne y$ in $V$ if $x+y$ is a square (forced pairs excluded)
\Procedure{Extend}{$\pi$} \Comment{$\pi$ a path in $V$ starting at $u$}
  \If{$|\pi|=|V|$} \Return $[\text{last}(\pi)=w]$ \EndIf
  \State $x\gets\text{last}(\pi)$
  \If{$x$ is an endpoint and its partner $x^*$ is unused}
     \State \Return \Call{Extend}{$\pi\cdot x^*$} \Comment{traverse the segment}
  \EndIf
  \For{each unused neighbour $y$ of $x$, with $y\ne w$ unless $|\pi|=|V|-1$}
     \If{\Call{Extend}{$\pi\cdot y$}} \Return true \EndIf
  \EndFor
  \State \Return false
\EndProcedure
\State \Call{Extend}{$(u)$} \Comment{exhaustive: complete search; randomised: shuffled neighbours, step budget, restarts}
\end{algorithmic}
\end{algorithm}

For $t=3$ the search is exhaustive and terminates without success in all $24$ cases. For
$t=5,7$ the randomised variant (step budget $2\cdot10^5$--$3\cdot10^5$, up to $60$ restarts)
found the arrangements of \cref{app:data-glue}. Each was then checked directly: every
junction sum is a square, and every segment occurs exactly once.

\subsection{Moves and certificates}

\begin{algorithm}[h]
\caption{Parity-preserving moves of a path $P=(v_1,\dots,v_N)$}\label{alg:moves}
\begin{algorithmic}[1]
\Function{Alternative}{$i,j$} \Comment{a parity-preserving re-ordering of window $[i,j]$, or none}
  \State $W\gets\{v_i,\dots,v_j\}$; $\mathrm{par}(v_p)\gets p\bmod 2$
  \State depth-first search for sequences $w_i,\dots,w_j$ of distinct elements of $W$ with
  \State \hspace{1.2em} $w_i\sim v_{i-1}$, $w_{p+1}\sim w_p$, $\mathrm{par}(w_p)=p\bmod 2$, $w_j\sim v_{j+1}$, and $(w_p)\ne(v_p)$
  \State \Return the first sequence found, or \textsc{none}
\EndFunction
\State $\mathcal W\gets\emptyset$
\For{$i=2,\dots,N-1$}
  \For{$j=i,\dots,\min(N-1,i+L-1)$} \Comment{$L$: maximum window length}
    \If{\Call{Alternative}{$i,j$}$\ne$\textsc{none}} add $[i,j]$ to $\mathcal W$; \textbf{break} \EndIf
  \EndFor
\EndFor
\State sort $\mathcal W$ by right endpoint; $\mathcal F\gets\emptyset$; $\ell\gets-\infty$
\For{$[i,j]\in\mathcal W$}
  \If{$i\ge\ell+2$} add $[i,j]$ to $\mathcal F$; $\ell\gets j$ \EndIf
\EndFor
\State \Return $\mathcal F$ \Comment{a maximum separated family among windows of length $\le L$}
\end{algorithmic}
\end{algorithm}

\begin{lemma}
\Cref{alg:moves} returns a separated family of maximum size among all families of windows of
length at most $L$ that admit moves.
\end{lemma}
\begin{proof}
Any window $[i,j']$ admitting a move contains the shortest such window $[i,j]$, $j\le j'$, which
starts at the same position. Replacing windows by these shorter ones preserves separation, so
it suffices to consider the shortest window at each start. Two windows are compatible if and only
if the closed intervals $[i,j+1]$ are disjoint, so this is interval scheduling, and the greedy
choice by right endpoint is optimal.
\end{proof}

\begin{algorithm}[h]
\caption{Verification of \cref{thm:main} for given seeds}\label{alg:certificate}
\begin{algorithmic}[1]
\Require a seed path $P$ of $G(n_0)$ from $1$ to $3$, $n_0$ odd; moves $(i_r,\omega_r)_{r\le k}$; the glue $\Gamma$ of \cref{tab:glue}; a number of levels $m$
\State check: $P$ is a permutation of $[n_0]$, consecutive sums are squares, $P_1=1$, $P_{n_0}=3$
\State check each move: $\omega_r$ re-orders the window, keeps parities, and all sums including those at the anchors are squares; check separation
\For{level $=1,\dots,m$}
  \State $(P',\mathrm{offset})\gets$ \Call{Lift}{$P,\Gamma$} (\cref{alg:lift}); check $P'$ as in line 1, with $n=25n_0+12$
  \For{each move $(i,\omega)$ of $P$, window $[i,j]$, and each $c\in[-12,12]$}
     \State $\omega^c_q\gets25\,\omega_q+(-1)^{i+q+1}c$ for $q=0,\dots,j-i$ \Comment{entry at position $i+q$ of $T_c$}
     \State $(o,\rho)\gets\mathrm{offset}[c]$ \Comment{$o$ = number of entries of $P'$ before the block}
     \State new move on window $[o+i,\,o+j]$ with content $\omega^c$ if $\rho$ is false,
     \State \hspace{1.2em} else on window $[o+n+1-j,\ o+n+1-i]$ with content $\omega^c$ reversed
  \EndFor
  \State check every new move as in line 2; check separation
  \State $P\gets P'$; $n\gets25n+12$; the moves $\gets$ the new moves
\EndFor
\end{algorithmic}
\end{algorithm}

\Cref{alg:certificate} performs only integer operations. We ran it with $m=2$ for the three
seeds of \cref{app:data-cert}. The results are in \cref{tab:liftcheck}.

\begin{algorithm}[h]
\caption{Seed search}\label{alg:seeds}
\begin{algorithmic}[1]
\Require odd $n$; a mode, exhaustive or random; a sample count $s$; a maximum window length $L$
\If{exhaustive}
  \State enumerate all Hamiltonian paths of $G(n)$ from $1$ to $3$ by depth-first search
\Else
  \State draw $s$ paths by randomised depth-first search: at each step try the unused neighbours in order
  \State \hspace{1.2em} of increasing number of unused neighbours (Warnsdorff's rule), ties broken at random,
  \State \hspace{1.2em} with a step budget of $2\cdot10^5$ and restarts
\EndIf
\State \Return $\max$ over the paths of $|\mathcal F|$ from \cref{alg:moves}
\end{algorithmic}
\end{algorithm}

\subsection{Numerical procedures}

\begin{algorithm}[h]
\caption{The constant $C$}\label{alg:sinkhorn-cont}
\begin{algorithmic}[1]
\Require grid size $M$, iteration count $I$
\State $x_a\gets(a-\frac12)/M$ for $a=1,\dots,M$; $K_{ab}\gets1/(2\sqrt{x_a+x_b})$; $\alpha\gets\mathbf1$
\For{$I$ iterations} $\alpha_a\gets\bigl(\alpha_a\big/\tfrac1M\sum_bK_{ab}\alpha_b\bigr)^{1/2}$ \EndFor
\State \Return $C\approx-\frac2M\sum_a\log\alpha_a$ and the residual $\max_a\bigl|\alpha_a\tfrac1M\sum_bK_{ab}\alpha_b-1\bigr|$
\end{algorithmic}
\end{algorithm}
We used $I=5000$ and $M\in\{400,1600,3200\}$. The midpoint rule handles the integrable corner
singularity adequately; the three values agree to $10^{-5}$.

\begin{algorithm}[h]
\caption{Sinkhorn scaling of $A_{G(n)}$, Bethe value, ansatz error}\label{alg:bethe}
\begin{algorithmic}[1]
\State build $A=A_{G(n)}$ in sparse form from $\mathrm{Adj}_n$; $x\gets\mathbf1$
\For{$3000$ iterations} $x\gets(x/Ax)^{1/2}$ \Comment{entrywise}\EndFor
\State $S\gets\operatorname{diag}(x)A\operatorname{diag}(x)$; check $\max_i|\sum_jS_{ij}-1|$ (below $2\cdot10^{-15}$ in all reported runs)
\State $F\gets\sum_{(i,j):A_{ij}=1}\bigl[-S_{ij}\log S_{ij}+(1-S_{ij})\log(1-S_{ij})\bigr]$ \Comment{\cref{cor:bethe}: $\log\operatorname{per}A\ge F$}
\State ansatz: $\widetilde p_{ij}=a(i/n)a(j/n)/\sqrt n$ with $a$ from \cref{alg:sinkhorn-cont}, linearly interpolated; report $\max_i$ and r.m.s.\ of $|\sum_j\widetilde p_{ij}-1|$
\end{algorithmic}
\end{algorithm}

\begin{algorithm}[h]
\caption{Finite Bregman bound}\label{alg:bregman}
\begin{algorithmic}[1]
\State $d_i\gets$ \eqref{eq:deg} for $i=1,\dots,n$
\State \Return $\log\mathrm{UB}(n)=\sum_i\frac{\log((d_i+1)!)}{d_i+1}+\frac{\log n!}{n}-\log2$ \Comment{$\Ham(n)\le\mathrm{UB}(n)$, \cref{prop:bregman}}
\end{algorithmic}
\end{algorithm}

\begin{algorithm}[h]
\caption{Edge count}\label{alg:edges}
\begin{algorithmic}[1]
\State \Return $e(G(n))=\sum_{k=2}^{\lfloor\sqrt{2n-1}\rfloor}\max\Bigl(0,\ \bigl\lfloor\tfrac{k^2-1}2\bigr\rfloor-\max(1,k^2-n)+1\Bigr)$
\end{algorithmic}
\end{algorithm}
The induced subgraph on the multiples of $4$ is isomorphic to $G(\lfloor n/4\rfloor)$ (\cref{prop:mod4}),
so its edge count is $e(G(\lfloor n/4\rfloor))$.

\begin{algorithm}[h]
\caption{Scaling predictor and residuals}\label{alg:predictor}
\begin{algorithmic}[1]
\For{each $n$ with $\Ham(n)>0$}
  \State $r_C\gets\log\Ham(n)-\bigl[\log(n!/2)-\frac{n-1}2\log n+Cn\bigr]$
  \If{$n\ge31$}
    \State dense Sinkhorn: $d\gets\mathbf1$; repeat $4\cdot10^4$ times $d\gets(d/Ad)^{1/2}$; $S\gets\operatorname{diag}(d)A\operatorname{diag}(d)$
    \State $\log\widehat\Ham(n)\gets\log\frac n2+\log n!-n\log n-2\sum_j\log d_j+2\log\bigl(\frac1n\sum_jd_j\bigr)$; $r_S\gets\log\Ham(n)-\log\widehat\Ham(n)$
    \State loop corrections, reported but not used: with $\lambda$ the eigenvalues of $S$ other than $1$,
    \State \hspace{1.2em} $-\sum_{\ell=1}^{n-1}(1-\ell/n)\bigl(\sum_\lambda\lambda^\ell\bigr)$ and $\sum_\lambda\log(1-\lambda)$
  \EndIf
\EndFor
\end{algorithmic}
\end{algorithm}
For $n\le30$ the iteration does not converge: some entries of $S$ tend to $0$ and the scaling
vector diverges, because $A_{G(n)}$ lacks total support. For $31\le n\le75$ the row sums of $S$
equal $1$ to within $5\cdot10^{-16}$.
